\documentclass[11pt]{article}
\usepackage[margin=1in]{geometry}
\usepackage{amsmath,amssymb,amsthm,lmodern}
\usepackage[hidelinks]{hyperref}
\hypersetup{pdftitle={Four interlacing rows of ordered set partitions at threshold three},
  pdfauthor={Bowen Zuo}}
\newtheorem{theorem}{Theorem}
\newtheorem{lemma}{Lemma}
\newcommand{\ii}{\mathrm{i}}
\newcommand{\sgn}{\operatorname{sgn}}
\newcommand{\Real}{\operatorname{Re}}
\newcommand{\Imag}{\operatorname{Im}}

\title{Four interlacing rows of ordered set partitions at threshold three}
\author{Bowen Zuo\\[2pt]\small East China University of Science and Technology}
\date{November 2026}

\begin{document}
\maketitle

\begin{abstract}
Let \(P_n(t)\) enumerate ordered set partitions of an \(n\)-element
labeled set by the number of blocks of size at least three. We prove
that every nonconstant row has only simple negative zeros and that
any four consecutive rows pairwise strictly interlace. This is the
largest window that holds uniformly in the starting rank. For
\(n\ge240\), a rectangular contour isolates one negative real pole
and a conjugate pair of poles in the generating function. The real
pole is eliminated exactly at each old root in the regime where all
three poles merge. Rational enclosures control the remaining
parameter range, and exact isolation certificates cover the 720
comparisons with smaller starting rank.
\end{abstract}

\section{The ordered model and the result}

An ordered set partition of \([n]\) is a sequence of nonempty,
pairwise disjoint subsets whose union is \([n]\). Every block takes
its place in that sequence, and the labels within a block are
unordered. Each block of size at least three contributes a factor
of \(t\). Write \(P_n(t)\) for the sum of these weights, including
\(P_0=1\) for the empty sequence. This is the threshold-three case
of Alexandersson's conjecture for the fully ordered model
\cite{symcat,realrooted}.

\begin{theorem}\label{thm:main}
For every \(n\ge0\), the polynomial \(P_n\) has degree
\(d=\lfloor n/3\rfloor\) and has \(d\) distinct negative roots. Write its
roots as \(a_1<\cdots<a_d\). For every \(h\in\{1,2,3\}\), let
\(b_1<\cdots<b_e\) be the roots of \(P_{n+h}\).
If \(e=d\) and \(d>0\), then
\[
 a_1<b_1<a_2<b_2<\cdots<a_d<b_d<0.
\]
If \(e=d+1\), then
\[
 b_1<a_1<b_2<a_2<\cdots<a_d<b_{d+1}<0.
\]
A constant has the empty root list. The second order in the case
\(d=0\) means \(b_1<0\), and two empty lists satisfy the first order
vacuously. Consequently any four consecutive rows pairwise strictly
interlace, and four is the largest such window valid for all starts.
\end{theorem}

To explain the relation to earlier models, set \(A(z)=z+z^2/2\) and
\(B(z)=e^z-1-z-z^2/2\). The exponential generating function is
\begin{equation}\label{eq:egf}
  \sum_{n\ge0}P_n(t)\frac{z^n}{n!}
  =\frac{1}{1-A(z)-tB(z)}.
\end{equation}
In the unordered model of the earlier question \cite{mo}, it is
\(e^{A(z)+tB(z)}\). The singleton-or-large lift described in the
catalog \cite{related} orders only the large blocks and excludes
pairs, giving \(e^z/(1-tB(z))\). Already at \(n=6\) these two models
give \(76+117t+10t^2\) and \(1+42t+20t^2\), respectively, whereas
\(P_6=3690+973t+20t^2\). The broader interlacing window was proposed
for the unordered problem in \cite{mo}. Fang \cite[Theorem~1]{fang}
proves eventual interlacing for the unordered threshold-four
polynomials and gaps one through four. Our conclusion concerns all
ranks at the fully ordered threshold three.

Selecting the label set of the last block gives the integer recurrence
\begin{equation}\label{eq:last}
 P_0=1,\qquad P_n(t)=\sum_{a=1}^{n}\binom na t^{[a\ge3]}P_{n-a}(t).
\end{equation}
In particular, \(P_1=1\), \(P_2=3\), and \(P_3=12+t\).
There are at most \(\lfloor n/3\rfloor\) marked blocks. For every
number of marked blocks within this bound, take that many triples and
fill the remaining labels with singleton blocks. Thus every
coefficient up to the degree is positive. A strictly alternating
pair of root lists has lengths differing by at most one. The degrees
of \(P_{3m+2}\) and \(P_{3m+6}\) differ by two, proving the asserted
obstruction to a uniform five-row window.

\section{The three poles}

Put \(t=-q/(1-q)\), with \(0<q<1\). This is a strictly decreasing
bijection onto the negative real axis. By \eqref{eq:egf},
\begin{equation}\label{eq:G}
 \frac{P_n(-q/(1-q))}{n!}=(1-q)[z^n]\frac{1}{G_q(z)},\qquad
 G_q(z)=1-2q-z-\frac{z^2}{2}+qe^z .
\end{equation}
At any zero \(z\) of \(G_q\), direct substitution yields
\begin{equation}\label{eq:deriv}
 G_q'(z)=\frac{z^2}{2}-2(1-q).
\end{equation}

Let \(J(x)=(1+x)e^{-x}-1+x^2/4\). Since
\(J'(x)=x(1/2-e^{-x})\), \(J\) has exactly one positive zero
\(\alpha\). Taylor bounds give \(J(1)<0<J(11/10)\).
Let \(q_c=1-\alpha^2/4=(1+\alpha)e^{-\alpha}\). Then
\(q_c>279/400\). For \(q_c<q<1\), the stationary equation
\(q=(1+x)e^{-x}\) has one solution in \((-1,0)\), a local maximum,
and one in \((0,\alpha)\), a local minimum. The minimum value is
\(-2J(x)>0\). Since \(G_q(0)>0\) and \(G_q(-2)<0\), there is exactly
one real zero, denoted \(-v\), with \(0<v<2\) and \(G_q'(-v)>0\).
On the negative axis to the left of the maximum \(G_q\) increases
strictly and to its right remains positive up to zero. The positive
minimum excludes all nonnegative real zeros.

\begin{lemma}\label{lem:rectangle}
Throughout \(q_c\le q\le1\), the boundary of
\(\mathcal R=[-2,2]+\ii[-\pi,\pi]\) satisfies
\(|G_q|\ge1/10\). There are precisely three zeros in its interior,
counted with multiplicity. For \(q_c<q<1\) they are the simple
real zero \(-v\) and a simple conjugate pair \(\zeta,\bar\zeta\)
with \(\Imag\zeta>0\). As \(q\downarrow q_c\), \(\zeta\to\alpha\).
\end{lemma}

\begin{proof}
Conjugation reduces the boundary checks to \(y=\Imag z\ge0\).
On the top edge the imaginary part is \(-(1+x)\pi\), where
\(x=\Real z\). When \(|x+1|<1/10\), the real part is at least
\(3/2-2+\pi^2/2-1/200-e^{-0.9}>1\). On the left edge the
imaginary part is at least \(y\), and for \(y<1/10\),
the real part is at most \(1-1.38+0.005+e^{-2}<-1/10\).
On the right edge the real and imaginary parts are
\[
  -3-2q+y^2/2+qe^2\cos y,\qquad qe^2\sin y-3y.
\]
For \(0\le y\le1/5\), use \(\cos y\ge0.98\) in the real part.
For \(1/5\le y\le1\) and \(1\le y\le\pi/2\), use
\(\sin y/y\ge5/6\) and \(\sin y/y\ge2/\pi\), respectively,
in the imaginary part. For \(\pi/2\le y\le2.3\) use
\(\cos y\le0\), and for \(2.3\le y\le\pi\) use
\(\cos y\le-1/2\), in the real part. These bounds use only \(q>0.69\)
and give \(|G_q|\ge1/10\).

At \(q=1\), Taylor's integral formula gives
\[
 G_1(z)=\frac{z^3}{2}\int_0^1(1-s)^2e^{sz}\,\mathrm ds .
\]
Multiplication by \(e^{-\ii\,\Imag(z)/2}\) makes the integral have
positive real part, since the rotated integrand has positive
real part for \(0<s<1\) and \(|\Imag z|\le\pi\).
The argument principle makes the count constant along \(q\).
For \(q_c<q<1\) there is one simple real zero. The other two are
nonreal conjugates and hence simple. The upper zero is continuous
in \(q\) by local simplicity and uniqueness. At \(q_c\), \(\alpha\)
is a double zero since \(G_{q_c}(\alpha)=G_{q_c}'(\alpha)=0\) and
\(G_{q_c}''(\alpha)=\alpha>0\); the negative zero remains simple.
A sufficiently small disk around \(\alpha\) retains two zeros by
Rouch\'e's theorem, giving the limit.
\end{proof}

On the open interval \(q_c<q<1\), write \(r=|\zeta|\),
\(\theta=\arg\zeta\), \(\beta=\arg G_q'(\zeta)\),
\(\lambda=r/v\), \(b=|G_q'(\zeta)|/(2G_q'(-v))\) and
\(\psi_n=(n+1)\theta+\beta\). The estimates below give
\(r<\min(2,v)\), \(0<\theta<\pi/3\) and \(0<\beta<\pi\).
Normalize with a positive factor and use the three residues:
\begin{align}
 S_n&=\tfrac12 |G_q'(\zeta)|r^{n+1}[z^n]G_q(z)^{-1}
      =-\cos\psi_n-b(-\lambda)^{n+1}+e_n,\label{eq:res}\\
 |e_n|&\le18|G_q'(\zeta)|(r/2)^{n+1}.\label{eq:error}
\end{align}
Indeed, the coefficient is the contour integral of
\(1/(G_q(z)z^{n+1})\) minus the other residues,
the integral being divided by \(2\pi\ii\). The boundary has length
\(8+4\pi\), with \(|z|\ge2\) and the lower bound in
Lemma~\ref{lem:rectangle}. The sign of \(S_n\) equals that of the
original row in \eqref{eq:G}.

\section{Rational bounds away from the merger}

Put \(\epsilon=(6(1-q))^{1/3}\). For \(q_c<q<1\),
\(0<\epsilon<61/50\). We first consider \(1/32\le\epsilon<61/50\).
The accompanying rational certificate partitions \([1/32,61/50]\)
into 851 closed intervals with shared endpoints. For each interval
the corresponding parameter bounds are
\(q_{\min}=1-\epsilon_{\max}^3/6\) and
\(q_{\max}=1-\epsilon_{\min}^3/6\).
Let \(q_0\) be their midpoint and \(\Delta q\) their half-width.
Each record supplies a disk with rational center \(c\), radius \(d\),
and two real rational numbers \(A_*>0,V_*>0\). The disk lies in
\(\Real z>0\), \(|\arg z|<\pi/3\), and \(|z|<A_*<V_*<2\).
The exact checks below are applied only to the oscillatory parameters
\(q_c<q<1\).

In 848 disks with \(\Imag c>d\), the comparison polynomial on
\(|w|=d\) is \(G_{q_0}'(c)w\). The residual has bound
\[
 |G_{q_0}(c)|+
 \tfrac12\bigl(|G_{q_0}''(c)|+q_0 e^{\Real c+d}d\bigr)d^2
 +\Delta q(2+e^{\Real c+d}).
\]
In the other three disks the center is real. The comparison polynomial
is \(G_{q_0}''(c)w^2/2\), and the residual has bound
\[
 |G_{q_0}(c)|+|G_{q_0}'(c)|d+
 \tfrac16 q_0e^{\Real c+d}d^3
 +\Delta q(2+e^{\Real c+d}).
\]
These follow from the Taylor formula on the complex line segment
and from \(G_q(z)-G_{q_0}(z)=(q-q_0)(e^z-2)\).
In each comparison the certificate strictly bounds the residual by
the modulus of the comparison polynomial. By Rouch\'e's theorem
the selected disk contains either the upper zero or the conjugate
pair, by Lemma~\ref{lem:rectangle} and the geometry of the disk.
The certificate checks \(G_q(-V_*)>0\) for both parameter endpoints,
which suffices by the affine dependence on \(q\).
The unique negative zero then gives \(V_*<v\). In particular
\(r<A_*<V_*<v\), and \(\Imag G_q'(\zeta)=\Real\zeta\,\Imag\zeta>0\)
by \eqref{eq:deriv}.

To see the error bound used by the checker, let \(X=\Real\zeta\),
\(Y=\Imag\zeta\). The imaginary equation at the root gives
\(qe^X\sin Y=(1+X)Y\), hence
\(\Real G_q'(\zeta)=(1+X)(Y\cot Y-1)\).
Since \(Y<\sqrt3\) and \(\sin Y\ge Y/2\) on this range,
\[
 1-Y\cot Y=\frac{\int_0^Y s\sin s\,\mathrm ds}{\sin Y}
       \le 2Y^2/3 .
\]
Using \(X<2\), \(r<2\), and \(\sin\theta=Y/r\), we obtain
\(|G_q'(\zeta)|/\sin\theta<12\). If \(\Imag c>d\), another
bound follows directly from \eqref{eq:deriv}. In that case the
certificate uses
\[
 g=\min\left(12,\frac{(A_*^2/2+2(1-q_{\min}))A_*}{\Imag c-d}\right);
\]
otherwise it takes \(g=12\).
The minimum of \(\sin(h\theta)/\sin\theta\), \(1\le h\le3\), is
\(\min(1,4\cos^2\theta-1)\). Since
\(4\cos^2\theta-1=(\sqrt3 X-Y)(\sqrt3 X+Y)/r^2\),
the certificate gives a positive lower bound \(H\) for that
minimum, using the two linear forms' Euclidean norm \(2\)
and the disk bounds on their values. Finally
\(G_q'(-v)>V_*^2/2-2(1-q_{\min})>0\).
Thus the checked inequality
\begin{equation}\label{eq:middle-cert}
 \frac{g}{H}\left(
   \frac{(A_*/V_*)^{241}}{2(V_*^2/2-2(1-q_{\min}))}
   +18(A_*/2)^{241}\right)<2/5
\end{equation}
gives, for \(m\ge240\) and \(h=1,2,3\),
\begin{equation}\label{eq:middle}
 b\lambda^{m+1}+|e_m|<(2/5)\sin(h\theta).
\end{equation}
Both bases in \eqref{eq:middle-cert} are below one, so
the exponent \(241\) suffices for the whole tail.
The largest upper bound rounded up to six decimals is \(0.385554\).
The verifier computes norm intervals by integer square roots and
uses the degree-36 Taylor polynomial with absolute remainder at most
\(8\cdot2^{37}/37!\) for all exponential inputs of modulus at most two.
In particular the intervals certify bounds throughout each cell in
the parameter and along each relevant disk boundary.

At an old root \(S_n=0\), \eqref{eq:middle} controls
\(|\cos\psi_n|\) and also the error in \(S_{n+h}\) by
\((2/5)\sin(h\theta)\). Expanding
\(-\cos\bigl(\psi_n+h\theta\bigr)\) in \eqref{eq:res},
the coefficient of \(\sin\psi_n\) is \(\sin(h\theta)>0\).
The two other terms together have modulus less than
\((4/5)\sin(h\theta)\), while
\(|\sin\psi_n|>\sqrt{1-(2/5)^2}>4/5\). It follows that
\begin{equation}\label{eq:sign}
 \sgn S_{n+h}=\sgn\sin\psi_n\quad\text{at every old root in this domain}.
\end{equation}

\section{Cancellation near the merger}

For \(0<\epsilon\le1/32\), scale \(z=\epsilon y\). Then
\[
 \frac{6G_q(\epsilon y)}{\epsilon^3}
 =2-e^{\epsilon y}
   +6\sum_{i\ge0}\frac{\epsilon^i y^{i+3}}{(i+3)!}
 =1+y^3+\epsilon(y^4/4-y)+O(\epsilon^2).
\]
For \(w^3=-1\), put \(a=5/12\) and \(c_w=w-a\epsilon w^2\).
The disk \(|y-c_w|<\epsilon^2\) contains one zero \(y_w\).
Indeed, on its boundary compare the scaled function with
\(3w^2(y-c_w)\). There \(|y-w|\le0.45\epsilon\) and
\(|y|\le1.02\). The first-order constant terms cancel
because \(w^4=-w\). The rest of the cubic part is at most
\(3.02|y-w|^2\), and the variation of \(y^4/4-y\) is
at most \((1+1.02^3)|y-w|\). The exponential series tails
are bounded using \(e^{1.02/32}<1.04\). Dividing the resulting
residual bound by \(\epsilon^2\) gives
\[
 3.02(0.45)^2+(1+1.02^3)0.45+
 1.04(1.02^2/2+1.02^5/20)<3 .
\]
Rouch\'e applies and the disks are disjoint. Thus \(w=-1\)
gives the real zero, and \(w=e^{\ii\pi/3}\) gives the upper zero.

Write \(u=r/\epsilon\), \(v_0=v/\epsilon\) and
\(\delta=\pi/3-\theta\). From
\(|1-a\epsilon w|^2=1-a\epsilon+a^2\epsilon^2\) at the upper
center, and the real center for \(w=-1\), we have
\[
 1-a\epsilon/2-\epsilon^2\le u
 \le1-a\epsilon/2+(1+a^2/2)\epsilon^2<1,\qquad
 |v_0-(1+a\epsilon)|\le\epsilon^2 .
\]
The upper disk center is below the ray of angle \(\pi/3\)
at perpendicular distance \(5\sqrt3\,\epsilon/24\).
The perturbation of this distance is at most \(\epsilon^2\).
Dividing by \(u\) and bounding the sine inverse gives
\[
 0.32\epsilon<\sin\delta<0.41\epsilon,\quad
 0<\delta<0.42\epsilon,\quad \sin(3\theta)>0.95\epsilon.
\]
The disk modulus bounds and \eqref{eq:deriv} imply
\[
 0.54\epsilon<1-\lambda<0.70\epsilon,\qquad
 b\le\frac{1+2\epsilon/3}{2(1-2\epsilon/3)}<0.522 .
\]
In particular \(r<v\) and \(r<\epsilon\).
Factoring \(G_q'(\zeta)=(\zeta^2/2)(1-4(1-q)/\zeta^2)\)
shows that the second factor differs from one by at most
\((2/3)\epsilon/u^2\). The disk bound on \(u\) gives
\(|\beta-2\theta|<0.72\epsilon\) and hence
\(|\beta-2\pi/3|<1.56\epsilon\).
Also \(|G_q'(\zeta)|<\epsilon^2\) by \eqref{eq:deriv}.
The scalar bounds just stated are checked by the accompanying
rational script. By \eqref{eq:error}, for \(m\ge240\) the
contour error is less than \(\epsilon/1000\).

Use these estimates in the portion of the search interval where
\begin{equation}\label{eq:n-epsilon}
 n\epsilon\ge3/2.
\end{equation}
There \(\lambda^{n+1}<e^{-0.81}<0.45\), so
\(|S_n+\cos\psi_n|<0.24\). At \(S_n=0\), subtract
\((-\lambda)^h S_n\) from \(S_{n+h}\) in \eqref{eq:res}. This removes
the real pole exactly and gives
\begin{equation}\label{eq:cancel}
 \begin{split}
 S_{n+h}={}&\sin\psi_n\sin(h\theta)
   +\bigl((-\lambda)^h-\cos(h\theta)\bigr)\cos\psi_n\\
 &\hspace{4mm}+e_{n+h}-(-\lambda)^he_n .
 \end{split}
\end{equation}
For \(h=1,2\), \(\sin(h\theta)>0.8\), and the bracket in
\eqref{eq:cancel} has modulus at most \(2\).
For \(h=3\), its modulus is at most
\(3(1-\lambda)+(3\delta)^2/2\), using
\(\cos(3\theta)=-\cos(3\delta)\). In all three cases the
bracket divided by \(\sin(h\theta)\) has modulus at most \(5/2\).
The two contour errors divided by that sine have sum of
moduli less than \(2/950\).
At an old root \(|\cos\psi_n|<0.24\), and
\(\sqrt{1-0.24^2}>(5/2)\cdot0.24+2/950\).
Therefore \eqref{eq:sign} also holds in the portion of the
merging regime specified by \eqref{eq:n-epsilon}.

\section{All roots and their order}

Fix \(n\ge240\) and set \(d=\lfloor n/3\rfloor\). We locate old
roots with
\[
 q_c<q\le q_*,
 \qquad q_*=1-\frac{1}{6}\left(\frac{3}{2n}\right)^3.
\]
On this interval \(n\epsilon\ge3/2\). The phase \(\psi_n\) is
continuous, since \(\zeta\) is simple and both arguments use the
sectors established above. At the left limit \(\theta\to0\) and
\(\beta\to\pi/2\). For the latter, as \(Y=\Imag\zeta\to0\),
the real part of \(G_q'(\zeta)\) is \(O(Y^2)\) by the formula
in Section~3 and its imaginary part is \(XY\) with \(X\to\alpha>0\).
At \(q_*\), the corner estimates give
\begin{align*}
 \left|\psi_n-(n+3)\pi/3\right|
 &<\bigl((n+1)0.42+1.56\bigr)\frac{3}{2n}\\
 &\le 0.642375<\pi/4 .
\end{align*}
Since \((n+3)\pi/3\ge(d+1)\pi\), the phase must traverse each strip
\[
 [i\pi+\pi/4,(i+1)\pi-\pi/4],\qquad 1\le i\le d.
\]
The assertion needs only continuity. For each strip take the first
hit of its upper endpoint after the preceding strip's upper hit,
and the last hit of its lower endpoint before that. Between these
two hits the phase lies within the strip, and the resulting
parameter intervals are disjoint and ordered.

On these intervals \(|S_n+\cos\psi_n|<2/5\) by \eqref{eq:middle}
and the bound \(0.24\) in the merging regime. The values of
\(-\cos\psi_n\) at the two boundary hits have opposite signs and
absolute value \(\sqrt{2}/2\). Hence each strip contains the phase
of a root at a parameter between its two hits. We have found
\(d\) distinct roots of the degree-\(d\) polynomial \(P_n\).
They are simple and account for all its roots. In particular,
there are no remaining roots in the region \(n\epsilon<3/2\).

At the \(i\)-th old root in increasing \(q\), the sign
of \(\sin\psi_n\) is \((-1)^i\). By \eqref{eq:sign},
this is the sign of the new row for each gap \(h=1,2,3\).
The new row is positive at \(q=0\) since its constant coefficient
is positive. Thus one new root lies before the first old root
in increasing \(q\), and one lies in each old gap. If the new
degree rises by one, its sign near \(q=1\) is \((-1)^{d+1}\),
because its leading coefficient is positive and \(t\to-\infty\).
This gives one more root after the last old root.
The new degree exhausts the roots in either case and makes each
simple. The map \(t=-q/(1-q)\) reverses the order, giving the two
orders of Theorem~\ref{thm:main}.

\section{The finite complement}

For \(0\le n\le239\) and gaps one through three, the exact certificate
uses the rows through 242. To verify it, reconstruct the coefficients
by cutting at marked blocks. Write \(V(z)=(1-A(z))^{-1}\).
A row with exactly \(k\) large blocks has a sequence of small
blocks in each of the \(k+1\) gaps. Hence
\[
 [t^k]P_n(t)=n![z^n]B(z)^k V(z)^{k+1}.
\]
An integer binomial convolution in exponential normalization computes
these coefficients. The generator proposes intervals by the
last-block recurrence \eqref{eq:last} and Bernstein subdivision.
The checker independently recomputes all rows by the marked-block
formula and uses integer Horner evaluation after multiplying by a
positive power of each endpoint denominator.

There are 9,720 isolating brackets across the 243 rows, disjoint
within each row and each with strict opposite endpoint signs.
Each row has as many brackets as its degree. Thus continuity and
degree give one simple negative root in every bracket and account
for all the roots. Within each pair of rows the complete combined
list of brackets has strict endpoint order, matching the two patterns
of Theorem~\ref{thm:main}. The checker verifies 237 pairs of equal
positive degree, 474 positive-degree pairs with a degree increase,
three constant--constant pairs and six constant--linear pairs.
This accounts for all 720 comparisons and completes the proof.

\section*{Certificates and formal verification}

The ancillary files contain the complete arithmetic certificates and
Python standard-library verifiers. The closed-interval tests, the
scalar corner checks and the finite-row comparisons were independently
replayed. The all-size paper proof also received an independent
mathematical review.
The Lean 4.33.1 formalization includes the all-size theorem
\path{OrderedPartitions.ThreePole.original_all}, its four-row form
\path{original_four_rows}, and the sharpness theorem
\path{five_rows_fail} in the same namespace. They refer to the original
polynomial and to ordered root arrays that exhaust the roots with
multiplicity. The formal proof connects the finite arithmetic
certificates and the full analytic argument to these declarations.
The accompanying Lean archive gives the pinned dependencies and
instructions for rebuilding and replaying the proof.

\paragraph{Disclosure.}
ChatGPT assisted with the analytic arguments, independent review,
certificate implementation, Lean work and manuscript preparation.

\begin{thebibliography}{9}
\bibitem{symcat}
P.~Alexandersson, \emph{Real-rooted polynomials: open problems},
the ordered set-partition conjecture, The Symmetric Functions Catalog.
\href{https://www.symmetricfunctions.com/realRootedOpenProblems.htm\#set-partitions-and-trees}{Catalog}.

\bibitem{realrooted}
P.~Alexandersson, \emph{[open problem] Which block-weight exponential
families are real-rooted?}, RealRooted issue \#2590, 18 November 2026.
The 21 November comment records finite evidence for the ordered family.
\url{https://github.com/PerAlexandersson/RealRooted/issues/2590}.

\bibitem{mo}
P.~Alexandersson, \emph{Set partitions with big blocks -- real-rooted
polynomials?}, MathOverflow question 462360, 17 January 2024.
\href{https://mathoverflow.net/questions/462360}{MathOverflow}.

\bibitem{related}
P.~Alexandersson, \emph{Real-rooted polynomials from set partitions},
The Symmetric Functions Catalog.
\href{https://www.symmetricfunctions.com/realRootedSetPartitions.htm}{Catalog}.

\bibitem{fang}
C.~Fang, \emph{Eventual interlacing of large-block set-partition
polynomials}, Hexagon preprint 2611.01129v1, 26 November 2026.
\href{https://hexagonmath.org/2611.01129v1}{Preprint}.
\end{thebibliography}
\end{document}
